% 12/04/2010 at 12:30
% Cinemática do Ponto Material: Movimento circular
% Written by TORDAR 

\documentclass[12pt]{article}  

\pagestyle{empty}

\usepackage{graphicx} 
\usepackage{pst-all}
\usepackage{pst-slpe}

%Let's go: 

\begin{document} 

\null\vskip4cm
\begin{center} 
\begin{pspicture}(-1.1,0)(1.1,1.1)
\psset{unit=8cm}
\psarc[linewidth=3pt,arcsepA=3pt,arcsepB=3pt,linestyle=dashed,dash=5mm 1mm,linecolor=gray](0,0){0.75}{0}{360}
\pscircle[linewidth=0.1mm,
linecolor=gray,
fillstyle=ccslope,
slopebegin=lightgray,
slopeend=black,
slopecenter=0.65 0.65,
slopesteps=25](0.5303,0.5303){0.05}
\psline[linewidth=3pt,linecolor=blue]{->}(-0.8,0)(0.85,0)
\psline[linewidth=3pt,linecolor=blue]{->}(0,-0.8)(0,0.85)
\psline[linewidth=5pt,linecolor=red]{->}(0,0)(0.5303,0.5303)
\rput{0}(0.9,0){\scalebox{1.5}{$X$}}
\rput{0}(0,0.9){\scalebox{1.5}{$Y$}}
%\rput{0}(0.3,0.15){\scalebox{1.5}{$\mathbf{r}(t)$}}
\end{pspicture}
\end{center}
\end{document} 
